%%%%%%%%%%%%
%
% $Autor: Wings $
% $Datum: 2019-03-05 08:03:15Z $
% $Pfad: SafetyGuidelines.tex $
% $Version: 4250 $
% !TeX spellcheck = en_GB/de_DE
% !TeX encoding = utf8
% !TeX root = manual 
% !TeX TXS-program:bibliography = txs:///biber
%
%%%%%%%%%%%%

\chapter{Allgemeine Hinweise}
Ein Umbau oder eine Erweiterung des Tischförderbands nach eigenem Ermessen des Benutzers ist nicht gestattet. Die Haftung seitens des Herstellers für dieses Produkts erlischt in dem Moment, in dem Änderungen an diesem Produkt vorgenommen werden, welche nicht dem Auslieferungszustand entsprechen. Selbiges gilt für Nichteinhaltung der folgenden Hinweise:

\section{Nutzungshinweise}
\begin{itemize}
	\item \textbf {Netzspannung:}	 
	Verwenden Sie das Netzteil nur mit der korrekten Netzspannung (100 – 240 V Wechselstrom). Schließen Sie es an eine gut erreichbare Steckdose in unmittelbarer Nähe des Tischförderbands an. Achten Sie darauf, dass das Netzkabel während der Benutzung nicht gestrafft ist. Achten Sie darüber hinaus darauf, dass das Netzkabel nicht in die Laufrollen oder das Transportband geraten kann.  
	\item \textbf {Oberfläche:}
	Stellen Sie das Tischförderband nur auf eine ebene, rutschfeste Tischfläche. Die Höhe der Tischfläche ist dabei so zu wählen, dass diese bei Benutzung des Tischförderbands im Stand nicht mehr als 10~cm über die Hüfthöhe der nutzenden Person hinausgeht.
	\item \textbf {Transportboxen:}
	Verwenden Sie ausschließlich die für das Tischförderband vorgesehenen und im Lieferumfang enthaltenen Transportboxen. Für andere Fördergüter ist dieses Tischförderband nicht ausgelegt.
	
\end{itemize}

\section{Aufbewahrung und Pflege}
\begin{itemize}
	\item \textbf {Reinigung:}
	Reinigen Sie das Produkt nur im stromlosen Zustand (Netzstecker ziehen). Wischen Sie es mit einem trockenen, weichen Tuch ab. Verwenden Sie keinesfalls Verdünner, Benzin oder aggressive Lösungsmittel, sowie Wasser. Ein eigenständiges Demontieren des Tischförderbands zu Reinigungszwecken ist nicht gestattet.
	
	\bigskip
	\bigskip
	\bigskip
	\bigskip
	\bigskip
	\bigskip
	\bigskip
	
	\item \textbf {Lagerung:}
	Setzen Sie das Tischförderband keinen hohen Temperaturen (durch Feuer oder andere Hitzequellen) oder direkter Sonneneinstrahlung aus. Lagern Sie das Tischförderband in trockener Umgebung bei Raumtemperatur. Lagern Sie das Tischförderband so ein, dass es mit eingeklappten Stützbeinen waagerecht. Es ist darauf zu achten, dass genügend Platz zu anderen Gegenständen gegeben ist (nicht unter 4cm). 

\end{itemize}

\section{Entsorgung}
\begin{itemize}
	\item \textbf {Entsorgung:} 
	Entsorgen Sie dieses Produkt nicht im Hausmüll. Da es sich bei dem Tischförderband um ein elektronisches Gerät handelt, ist dieses an entsprechender Stelle (z.B. kommunaler Wertstoffhof) zu entsorgen. Demontieren Sie das Tischförderband nicht eigenhändig, um eine Wertstofftrennung vorzunehmen.
	\item \textbf {Rückgabe:}
	In Deutschland sind Vertreiber von Elektrogeräten verpflichtet, Altgeräte unter bestimmten Bedingungen unentgeltlich zurückzunehmen. Adressen kostenfreier Sammelstellen erhalten Sie bei Ihrer Kommunalverwaltung.
	\item \textbf {Datenschutz:}
	Sollte das Gerät personenbezogene Daten enthalten, sind Sie selbst für deren Löschung verantwortlich.
\end{itemize}